% =====================================================================
%  Wspólna preambuła szablonów pism ZUS (świadczenia rodzicielskie 2026)
%  Silnik: XeLaTeX albo Tectonic (pdflatex NIE obsłuży polskich znaków w fontspec).
%  Kompilacja:  tectonic -X compile plik.tex
% =====================================================================
\documentclass[11pt,a4paper]{article}

\usepackage{fontspec}
\usepackage{polyglossia}
\setmainlanguage{polish}
\setotherlanguage{english}
% font domyślny (Latin Modern) — Tectonic go dostarcza

\usepackage[margin=2.5cm,a4paper]{geometry}
\usepackage{graphicx}
\usepackage{xcolor}
\usepackage{array}
\usepackage{enumitem}
\usepackage{fancyhdr}
\usepackage{hyperref}

\definecolor{zus}{HTML}{007A33}      % zieleń ZUS
\definecolor{szary}{HTML}{6B7280}
\definecolor{ramka}{HTML}{9CA3AF}

\hypersetup{colorlinks=true, linkcolor=zus, urlcolor=zus, pdftitle={Szablony pism ZUS}}

% --- pole do wypełnienia odręcznego -----------------------------------
% \Pole{nazwa} rysuje kropkowaną linię z podpisem pola pod spodem.
% Pole do wypełnienia: kropkowana linia, pod nią drobny podpis pola.
\newcommand{\Pole}[1]{%
  \textcolor{szary}{\dotfill}\,%
  {\scriptsize\textcolor{szary}{\textsc{#1}}}%
}

% --- blok warunkowy: treść pokazywana tylko, gdy pole prawdziwe --------
\newenvironment{warunek}{%
  \par\smallskip\noindent\begin{minipage}{\linewidth}%
  \color{black}%
}{%
  \end{minipage}\par\smallskip%
}

% --- nagłówek pisma ---------------------------------------------------
\newcommand{\NaglowekPisma}[1]{%
  \par\vspace{0.8em}%
  \begin{center}%
    {\large\bfseries\MakeUppercase{#1}}%
  \end{center}%
  \vspace{0.4em}%
}

% --- blok nadawcy (pierwsze linie pisma) -------------------------------
\newenvironment{nadawca}{%
  \begin{flushleft}\small
}{%
  \end{flushleft}\normalsize
}

% --- stopka z numerem strony -----------------------------------------
\pagestyle{fancy}
\fancyhf{}
\renewcommand{\headrulewidth}{0pt}
\fancyfoot[C]{\textcolor{szary}{\small\thepage}}

% --- metadane szablonu na początku pisma ------------------------------
\newcommand{\Metadane}[3]{%
  \begin{flushright}\footnotesize\textcolor{szary}%
    \textbf{#1}\\[2pt] {\itshape #2}\\[2pt] #3%
  \end{flushright}\vspace{0.5em}%
}

\setlength{\parindent}{0pt}
\setlength{\parskip}{0.6em}

\begin{document}

\begin{flushright}
{\footnotesize\color{szary}
\textbf{Szablon P01}\\
{\itshape Wniosek o udzielenie urlopu macierzyńskiego}\\
art. 180 \S{} 1 k.p. -- Formularz: ZAM}\
\end{flushright}

\vspace{0.5em}
{\color{ramka}\hrule height 0.4pt}
\vspace{0.8em}

\Pole{imie}\\[-2pt]
\Pole{adres}\\[-2pt]
PESEL: \Pole{pesel}

\Pole{platnik nazwa}
\Pole{platnik adres}
NIP: \Pole{platnik nip}

\Pole{data pisma}, \Pole{adres miejscowosc}

\NaglowekPisma{WNIOSEK O UDZIELENIE URLOPU MACIERZYŃSKIEGO}

Zgodnie z art. 180 \S{} 1 Kodeksu pracy wnoszę o udzielenie mi urlopu
macierzyńskiego w wymiarze \Pole{tygodnie} tygodni, przysługującego w związku
z urodzeniem \Pole{liczba dzieci pl} dziecka (dzieci) przy jednym porodzie.

Poród nastąpił dnia \Pole{data porodu}.
Urlop macierzyński ma trwać od dnia \Pole{urlop od} do dnia \Pole{urlop do}.

\begin{warunek}
Dziecko: \Pole{dziecko imie}, PESEL \Pole{dziecko pesel}.
\end{warunek}

\begin{warunek}
Do wniosku załączam zaświadczenie szpitalne, że dziecko urodziło się martwe.
\end{warunek}

\begin{warunek}
Jednocześnie, na podstawie art. 179\textsuperscript{1} Kodeksu pracy, wnoszę o udzielenie mi\\[-2pt]
bezpośrednio po urlopie macierzyńskim urlopu rodzicielskiego w pełnym wymiarze.\\[-2pt]
Oświadczam, że zasiłek macierzyński za cały okres będzie wynosił 81,5\% podstawy.
\end{warunek}

\Pole{imie}
......................................
(podpis pracownika)

\vfill
{\footnotesize\textcolor{szary}\itshape Przed porodem wpisuje się przewidywaną datę porodu i uzupełnia po porodzie. Urlopy macierzyński i rodzicielski łącznie liczą się do limitu art. 180 \S{} 1 k.p. }

\end{document}
